% ============================================================
%  DESIGN CONCEPT SUMMARY
% ============================================================

The LBR is a heavy-water-moderated, heavy-water-cooled reactor operating at
10\,MPa with a thermal output of 2.0\,GW$_\text{th}$. Spent PWR fuel
assemblies are accepted without mechanical or chemical processing; each
individual fuel rod is enclosed in a close-fitting Zircaloy-4 overclad sleeve
whose end-caps are remote-welded, helium-backfilled, and leak-tested before
integration into the core structure. The design philosophy is one of minimum
intervention: every refabrication step that can be eliminated is eliminated.

% ============================================================
\subsection{Core Layout and Operating Parameters}
% ============================================================

The active core is composed of assemblies arranged on a 6\,cm square lattice
pitch. Each assembly is built around a 15$\times$15 structural skeleton;
of the 225 available lattice positions, 192 are occupied by overcladded spent
fuel rods and 25 are occupied by Fueled Power Inverter (PI-F) devices
distributed at alternating positions throughout the assembly. At a specific
power of 30\,W/gHM, the 2\,GW$_\text{th}$ thermal output requires
approximately 67\,tonnes of heavy metal distributed across $\sim$19\,400 fuel
rods. Primary operating parameters are summarized in
Table~\ref{tab:core_params}.

\begin{table}[H]
\centering
\caption{Summary of primary LBR core design parameters.}
\label{tab:core_params}
\renewcommand{\arraystretch}{1.35}
\begin{tabular}{l l}
\toprule
\textbf{Parameter} & \textbf{Value} \\
\midrule
Thermal output              & 2.0\,GW$_\text{th}$ \\
Moderator / coolant         & Heavy water (D$_2$O) \\
Operating pressure          & 10\,MPa \\
Coolant inlet / outlet      & 533\,K / 582\,K \\
Specific power              & 30\,W/gHM \\
Lattice pitch               & 6\,cm (square) \\
Assembly layout             & 15$\times$15 \\
Spent fuel positions / assembly & 192 \\
PI-F positions / assembly   & 25 \\
Total heavy metal inventory & $\approx$67\,t \\
Total fuel rods (core)      & $\approx$19\,400 \\
Design discharge burnup     & 60.1\,GWd/tHM \\
Power peaking factor (PPF)  & 2.23 \\
DNBR (hot channel)          & 2.95 \\
\bottomrule
\end{tabular}
\end{table}

The spent fuel feedstock is characterized by an initial enrichment of
$\sim$4.01--4.5\,wt.\% $^{235}$U and a discharge burnup of
$\sim$47--53\,GWd/tHM from PWR service, representative of assemblies from
the U.S. commercial fleet. Fuel compositions were generated by a
Serpent~2 depletion model benchmarked against the TMI-1 NJ05YU assembly
post-irradiation examination data and transferred to the LBR lattice model
via binary restart file, preserving the full axially resolved nuclide
inventory.

% ============================================================
\subsection{Overclad Sleeve and Fuel Rod Preparation}
% ============================================================

The overclad is a close-fitting Zircaloy-4 tube into which a spent PWR fuel
rod is inserted without breaching the original cladding. End-caps are
remote-welded under hot-cell conditions, after which the annular space is
helium-backfilled to a specified pressure, leak-tested, and subjected to
non-destructive assay to confirm fuel rod identity. The overclad outer
diameter and wall thickness are chosen to maintain the same rod outer
dimension as the PI-F rods, so that all lattice positions share a uniform
pitch geometry. This design preserves the original spent fuel cladding as a
first barrier, adds the overclad as a second barrier, and avoids all
high-activity powder handling operations. The result is a dramatically
simplified hot-cell processing chain relative to DUPIC, as quantified in
Section~\ref{sec:economics}.

% ============================================================
\subsection{Fueled Power Inverter Device}
% ============================================================

The Fueled Power Inverter is the primary novel component of the LBR assembly.
It is a fresh UO$_2$ fuel rod with a graded axial enrichment profile designed
to concentrate fission power in the axially under-depleted extremities of
adjacent spent rods---the regions where PWR irradiation history leaves the
highest residual fissile inventory.

The device geometry matches the outer dimensions of the overcladded spent fuel
rods: Zircaloy-4 cladding with a 10.92\,mm outer diameter and an active fuel
length of 360.17\,cm. The axial enrichment profile consists of two zones:
\begin{itemize}
    \item \textbf{End regions} (upper and lower $\sim$70\,cm): 4.5\,wt.\%
    $^{235}$U --- co-located with the under-depleted extremities of the
    surrounding spent rods, where the additional reactivity drives enhanced
    local fission rates and ``inverts'' the axial power peak toward the rod
    ends.
    \item \textbf{Central region} ($\sim$220\,cm): 3.5\,wt.\% $^{235}$U ---
    co-located with the high-burnup midplane of adjacent spent rods, where
    lower enrichment limits the power contribution from the fresh fuel
    without suppressing criticality.
\end{itemize}

This configuration was selected from a parametric enrichment sweep (central
enrichments 1.25--8.25\,wt.\%, edge enrichments always $>$ central) on the
basis of a weighted decision matrix prioritizing critical cycle length
(weight 0.40) and power peaking factor (weight 0.35). The 3.5\,/\,4.5\,wt.\%
combination uniquely satisfies both the 60\,GWd/t burnup objective and the
PPF\,$\leq 2.5$ safety objective. The resulting bimodal axial power
distribution---fission peaks at the rod ends flanking a lower-power
midplane---is the defining neutronics signature of the LBR design. A detailed
description of Power Inverter geometry, materials, and fabrication
considerations is provided in Appendix~\ref{sec:app_A}.

% ============================================================
\subsection{Relationship to CANDU and the DUPIC Fuel Cycle}
% ============================================================

The LBR shares architectural heritage with the CANDU reactor family. Both
systems exploit D$_2$O moderation to maintain criticality on sub-enriched
fuel, and both can in principle be operated with online or near-online
refueling that eliminates the large initial excess reactivity burden of batch
refueling. The closest antecedent to the LBR concept is the DUPIC (Direct Use
of PWR Fuel In CANDU) cycle, which also proposes loading spent PWR fuel into
a heavy-water reactor without chemical separation. Table~\ref{tab:dupic_lbr}
summarizes the key distinctions.

\begin{table}[H]
\centering
\caption{Comparison of DUPIC and LBR fuel cycle process steps and key
metrics (March~2026 USD).}
\label{tab:dupic_lbr}
\renewcommand{\arraystretch}{1.30}
\begin{tabular}{l c c}
\toprule
\textbf{Attribute} & \textbf{DUPIC} & \textbf{LBR} \\
\midrule
Decladding required             & Yes               & No \\
OREOX cycling                   & Yes               & No \\
High-temperature sintering      & Yes (1650--1850\,$^\circ$C) & No \\
Pellet grinding                 & Yes               & No \\
New cladding required           & Yes (CANFLEX)     & Overclad sleeve only \\
Assembly geometry               & 43-element CANFLEX & 15$\times$15, 6\,cm pitch \\
Cladding-hull waste stream      & Yes               & \textbf{Eliminated} \\
Fabrication cost (\$/kgHM)      & 1,195             & 501 \\
All-in lifecycle cost (\$/kgHM) & 2,353             & 1,659 \\
Incremental burnup (MWd/kgHM)   & 13                & 15 \\
Incremental energy cost (\$/MWd)& 181               & 111 \\
\bottomrule
\end{tabular}
\end{table}

The key differentiator is the elimination of the thermo-chemical refabrication
chain. Because the original cladding is never breached, LBR does not generate
the Zircaloy cladding-hull waste stream (0.29\,t/tHM in DUPIC) that
constitutes the dominant component of DUPIC's process radwaste cost, and it
avoids the OREOX cycling, powder handling, and sintering operations that
account for $\sim$65\% of DUPIC's direct capital cost. The tradeoff is
reduced geometric flexibility: the LBR must accommodate the existing PWR rod
dimensions and axial composition profile, which drove the development of the
Power Inverter concept as a power-shaping tool.

% ============================================================
\subsection{Materials Summary}
% ============================================================

The primary materials employed in the final LBR design are:

\begin{itemize}
    \item \textbf{Fuel}: Commercial UO$_2$ as discharged from PWR service,
    with initial enrichment $\approx$4.01--4.5\,wt.\% and discharge burnup
    $\approx$47--53\,GWd/t. The PI-F rods use fresh UO$_2$ at graded
    enrichments of 3.5\,wt.\% (center) and 4.5\,wt.\% (edge).

    \item \textbf{Cladding and overclad}: Zircaloy-4, selected for its
    established neutron irradiation performance database, compatibility with
    pressurized D$_2$O coolant at 10\,MPa and 533--582\,K, and low thermal
    neutron absorption cross section. Thermal conductivity is conservatively
    taken as 9.2\,W/m-K (degraded spent-fuel value) for thermal analysis.

    \item \textbf{Moderator and coolant}: Heavy water (D$_2$O) at 10\,MPa.
    Properties used in thermal-hydraulic analysis: $\rho = 874$\,kg/m$^3$,
    $c_p = 4670$\,J/kg-K, $k = 0.53$\,W/m-K, $\mu = 1.14\times10^{-4}$\,kg/m-s.

    \item \textbf{Structural components}: Stainless steel (SS-304) for
    assembly grid spacers, end fittings, upper and lower plenums, and pressure
    boundary hardware, consistent with the TMI-1 reference assembly geometry.
\end{itemize}

The structural integrity and lifetime performance of the fuel-cladding system
under irradiation, thermal cycling, and corrosion were evaluated in
Section~\ref{sec:Materials} using the MOOSE multiphysics finite element
framework, coupling heat conduction, solid mechanics, creep deformation, and
corrosion-induced property degradation.
